\documentclass[12pt]{article}
\usepackage[utf8]{inputenc}
\usepackage[T1]{fontenc}
\usepackage{hyperref}
\title{Nisu, Northern Language Resources for K-12 Curriculum Development}
\author{Compiled by Claude Sonnet 5.5 (claude-sonnet-5-5) for Andr\'as Kornai}
\date{}
\begin{document}
\maketitle

\section{Introduction}
Northern Nisu (ISO 639-3 yiv; Glottocode nort2718) is a Loloish language of Yunnan, China. It is distinct from Nuosu (``Northern Yi''), which has a standard script and school system. Nisu varieties carry several ISO codes: nsd (Southern), yiv (Northern), nos (Eastern), nsv (Southwestern) and nsf (Northwestern). Little specific to Northern Nisu turned up in the searches.

\section{Primary Educational Resources}
No primers, textbooks or school programmes for Northern Nisu were found.

\section{Linguistic and Reference Materials}
\begin{itemize}
\item Wikipedia, Nisu language (dialect classification, citing Yang 2009 and Chen et al. 1985): \url{https://en.wikipedia.org/wiki/Nisu_language}
\item Wikipedia, Northern Loloish languages: \url{https://en.wikipedia.org/wiki/Northern_Loloish_languages}
\item Wikipedia, Nuosu language (a different language, listed to avoid confusion): \url{https://en.wikipedia.org/wiki/Nuosu_language}
\item Glottolog, Southern Nisu (nsd), a different variety: \url{https://glottolog.org/resource/languoid/iso/nsd}
\end{itemize}
A 2021 Nisu--Chinese dictionary (IPA and Yi script) was mentioned in a search summary. It is based on the Nisu of villages north of Mengzi, and I did not confirm that it covers Northern Nisu. No URL for it was among the results, so it is not listed as a resource.

\section{Religious and Community Resources}
None found.

\section{Digital and Audio Resources}
None found.

\section{Curriculum Development Context}
The language is used by all adults but not all young people, and it is not known to be taught in schools. Existing Yi-script and Nuosu-based school materials belong to another language and should not be assumed to fit Northern Nisu.
\end{document}